\chapter{Experiments and Results}
\label{chap:results}

All Approach~C experiments use \texttt{qwen3.6-plus} on the $13$-clip pen-conveyor
dataset under the blind protocol of Section~\ref{sec:vera}. Approach~A and B use the
R01 push-button scene.

\section{Approach A: Direct Prompting}
The direct-prompting exploration confirmed two qualitative findings that motivated the
later work: (i) step-wise / chain-of-thought prompts give more consistent verdicts than
terse single-paragraph prompts, and (ii) under compact output formats the model is
biased toward ``normal'', producing false negatives. These observations directly
motivated the articulate-first format (Approach~B) and the guiding-question
decomposition (Approach~C).

\section{Approach B: Multi-Persona Voting}
A single-prompt VLM verdict is unreliable on edge cases, so this approach adopts a
\emph{self-consistency} / persona-prompting paradigm: the same clip is inspected by
several diverse ``inspector personas'' and the decision is taken by majority vote. This
section details the system, its three prompt-design variants, and the qualitative
findings, in the order they were developed.

\subsection{Local Deployment}
Unlike the API-served models of Approaches~A and~C, the persona pipeline runs a locally
deployed open-weight model, \texttt{Qwen3-VL-8B-Instruct} \cite{qwen3}, through the HuggingFace
\texttt{transformers} stack. Running on-premises means no video ever leaves the factory,
which is a practical requirement for industrial data, and it also gives full control over
decoding: we use greedy decoding (\texttt{do\_sample=False}), so each persona is
deterministic and the only source of diversity is the prompt itself. Video is fed at
\texttt{fps}$=2$, capped at $32$ frames and $360\times420$ pixels, with up to $1024$ new
tokens per persona.

\subsection{The Twenty Personas}
Twenty distinct ``inspector personas'' were authored, grouped as four production-line
operators (senior, new, supervisor, night-shift), five engineers (mechanical, quality/Six
Sigma, robotics, computer-vision, industrial), and eleven further profiles (maintenance
technician, internal auditor, lean specialist, EHS specialist, academic, Toyota expert,
supplier-quality engineer, trainer, conservative auditor, aggressive auditor, hybrid).
Crucially, every persona is given the \emph{same} anomaly rules; they differ only in
voice, risk threshold and output format (for example the conservative auditor only raises
an alarm on clear multi-frame evidence, while the aggressive auditor flags any hint of a
fault). This yields \emph{synthetic diversity} for ensembling without changing the
underlying task.

\subsection{Voting Mechanism}
From the twenty personas, a subset of five (one operator, one mechanical engineer, one
vision engineer, one Toyota expert, one hybrid) plus a distilled \emph{consensus} prompt
form six voters. Each voter's free-text output is parsed to NORMAL / ANOMALY / UNCLEAR by
matching its persona-specific decision markers (e.g.\ \texttt{RESULT}, \texttt{PART
STATUS}, \texttt{ANDON CALL}, \texttt{ROBOTIC FITNESS}), with care taken not to misread
``no anomaly''. The final verdict is the majority over the decisive (non-UNCLEAR) votes,
and the confidence is the winning fraction ($\geq 0.8$ HIGH, $\geq 0.6$ MEDIUM, else LOW).
The consensus prompt is itself a contribution: it is a single prompt distilled from the
common ground of all twenty personas, used as a strong stand-alone voter.

\subsection{The Three Ablation Variants}
The personas were studied at three checklist breadths (defined in
Table~\ref{tab:persona-types}):
\begin{itemize}
  \item \textbf{specialized} --- each voter inspects a single focus area, so the committee
        divides the labour across orientation, motion and belt-position specialists.
  \item \textbf{complete} --- every voter runs the full A/B/C checklist (orientation,
        motion/alignment, belt position), providing redundancy: several independent full
        inspections are combined.
  \item \textbf{orientation} --- the multi-check structure is dropped entirely and every
        voter answers only the primary question for this line, ``is the red cap on the
        wrong side?''. This isolates the dominant anomaly type.
\end{itemize}
The variants form an ablation on \emph{how much structure each voter needs}: a minimal
single question, a single specialised focus, or the full checklist.
Table~\ref{tab:persona-proscons} summarises their trade-offs.

\begin{table}[!htbp]
\centering
\footnotesize
\caption{\label{tab:persona-proscons}Trade-offs of the three persona prompt designs.}
\begin{tabular}{l|p{3.5cm}|p{4.6cm}}
\hline
Design & Pros & Cons \\
\hline
specialized      & deep per-aspect reasoning; lower prompt overhead
                 & systematic blind spots; voters miss anomalies outside their domain \\
complete         & no blind spots; multi-axis verification
                 & verbose prompts; slower inference; bias toward compact YES/NO \\
orientation-only & sharpest signal; fastest inference; best alignment with the dominant
                   defect
                 & cannot detect motion or position anomalies \\
\hline
\end{tabular}
\end{table}

\subsection{The Articulate-First Finding}
The central qualitative result emerged while tuning the output format. With compact
YES/NO answer formats the model produced systematic \emph{false negatives}: it defaulted
to ``normal'' without actually examining the frames. Forcing an \emph{articulate-first}
format --- the voter must first describe the entry, middle and exit frames, then derive
each check, then state the verdict --- suppressed this bias and aligned the persona
verdicts with the (correct) consensus prompt. This is the same principle that underlies
Approach~C's guiding questions: compelling the model to verbalise its observations before
committing to a label is what improves reliability, independent of the specific wording.

\subsection{Status and Limitation}
The persona pipeline was developed on the R01 push-button scene, whose dominant anomaly is
the reversed orientation. A full quantitative evaluation of the three ablation variants on
a labelled multi-clip set was \emph{not} completed: the only saved voting report is a
single demonstration clip, for which all six voters returned UNCLEAR. The contributions of
this approach are therefore the system design (twenty personas, the consensus prompt, the
six-voter majority scheme), the three checklist-breadth ablations, and the articulate-first
finding; a fair numerical comparison of the variants, ideally on the same labelled set as
Approach~C, is left as future work (Chapter~\ref{chap:conclusion}).

\section{Approach C: VERA-lite}
This section describes, in the order they were tried, every method we evaluated within
the VERA-lite pipeline, and then compares them. Each step was motivated by the failure of
the previous one.

\subsection{Perception Probing with a Describe Prompt}
Before writing any classifier prompt, we ran a neutral \emph{describe} prompt that asked
the model only to report what it saw (object count, colour, orientation, motion
continuity), explicitly forbidding a verdict. This revealed how the model actually
perceives the scene: it counts the three pens reliably, sees the rail-and-board setup,
but perceives colour weakly (``greyish'' rather than ``blue'') and is unsure whether the
dark end is a cap. These observations shaped the guiding questions toward
\emph{relative} comparisons among the three pens (``are all three the same colour?'')
rather than absolute attributes.

\subsection{Blind Evaluation and Label-Leakage Correction}
The original pipeline passed the expected label into the learner prompt. A first run that
scored $13/13$ was therefore untrustworthy. After making blind mode the default, the
first honest score dropped to $0.77$; all subsequent numbers are blind.

\subsection{Guiding-Question Iteration (v2 $\rightarrow$ v3)}
For every clip the learner returns a structured JSON object: for \emph{each} guiding
question a YES/NO/UNCLEAR answer with a short visual justification, followed by the overall
verdict (NORMAL/ANOMALY), an anomaly type, an anomaly score and a confidence. A concrete
instance of this output is shown later in Table~\ref{tab:example}; this per-question
evidence is what makes the verdict explainable.

The describe step above gave us a first, hand-written question set, \textbf{v2}. We then
ran one VERA-style optimiser iteration: the optimiser VLM read the learner's answers and
rewrote the questions. Its verbatim rationale was that the v2 questions ``rely heavily on
negative phrasing $(\ldots)$ [and] VLMs frequently struggle with logical negation,'' so it
converted them into ``direct, positive verification statements'' with one atomic check per
question. On top of the optimiser's output we hardened the orientation question
(``including a pen near the frame edge''), the motion question (``including a brief
$1$--$3$\,s pause''), and added a sixth question for a pen moving relative to the board,
yielding \textbf{v3}. Table~\ref{tab:questions} places the two question sets side by side,
and Table~\ref{tab:qiter} reports the resulting accuracy: $0.77 \rightarrow 0.92$ (single
run), with the only remaining miss being the intra-object temporal anomaly.

\begin{table}[!htbp]
\centering
\footnotesize
\caption{\label{tab:questions}The guiding questions before (v2) and after (v3) the
optimiser iteration, by anomaly family. v2 uses negated phrasing; v3 is positive with one
atomic check per question, adds an edge-aware orientation check and a brief-pause motion
check, and introduces a sixth pen-movement question.}
\begin{tabular}{l|p{4.4cm}|p{5.3cm}}
\hline
Anomaly family & v2 (negated) & v3 (positive / hardened) \\
\hline
Count        & no missing pen and no empty gap      & exactly three pens in a single row \\
Orientation  & none flipped or reversed             & all pointing the same way, including a pen near the frame edge \\
Colour       & none differing in colour             & identical barrel and cap colours, no mismatch \\
Completeness & no bare tip, none shorter            & fully capped, no exposed tip, all equal length \\
Belt motion  & without stopping, pausing, reversing & advances at every moment, never halting even briefly ($1$--$3$\,s) \\
Pen movement & (not asked)                          & pens stay fixed; no pen sliding relative to the board \\
\hline
\end{tabular}
\end{table}

\begin{table}[!htbp]
\centering
\caption{\label{tab:qiter}Guiding-question iteration (blind, single run).}
\begin{tabular}{l|c|c|l}
\hline
Question set & Accuracy & FP & Missed clips \\
\hline
v2 & 0.77 (10/13) & 0 & orientation-2, stop-2, temporal \\
v3 & 0.92 (12/13) & 0 & temporal \\
\hline
\end{tabular}
\end{table}

\subsection{A Worked Example}
To illustrate the explainability of the approach, Table~\ref{tab:example} shows the
learner's actual (blind) output on the \texttt{wrong\_colour} clip. The model answers
each guiding question with a YES/NO and a short visual justification; the single NO on the
colour question (Q3) drives the final verdict. The structured per-question evidence lets
an operator see exactly \emph{why} the clip was flagged.

\begin{table}[!htbp]
\centering
\caption{\label{tab:example}Learner output on the \texttt{wrong\_colour} clip (verdict:
ANOMALY, type: wrong\_colour, score: $1.0$, confidence: HIGH).}
\begin{tabular}{c|c|p{9cm}}
\hline
Q & Ans. & Evidence (model's own words) \\
\hline
1 & YES & At 00:10 the board is fully visible and holds three pens in a row. \\
2 & YES & All three pens are oriented vertically with caps pointing down. \\
3 & \textbf{NO} & The leftmost pen has a blue barrel, while the middle and right pens are black/dark grey. \\
4 & YES & All pens appear capped and of similar length. \\
5 & YES & The board moves continuously left to right without stopping. \\
6 & YES & The pens remain fixed in their relative positions as the board slides. \\
\hline
\end{tabular}
\\[3pt]
\footnotesize{Reason: ``three pens on the board, but one is blue while the other two are
black, violating the requirement that all pens be identical in colour.''}
\end{table}

\subsection{Per-Clip Output of the v3 Run}
Table~\ref{tab:perclip} lists the learner's verdict, predicted anomaly type and score for
all $13$ clips in the $0.92$ (v3, single-run) configuration; the raw JSON output for one
clip is reproduced in Appendix~A. The model reported HIGH confidence on every
clip. On the combined clip it correctly returns ANOMALY but names only one of the two
co-occurring anomaly types, and the temporal clip is the single, confidently wrong, miss.

\begin{table}[!htbp]
\centering
\footnotesize
\caption{\label{tab:perclip}Per-clip output of the v3 video run (blind, single run;
confidence was HIGH on every clip).}
\begin{tabular}{l|c|c|l|c|c}
\hline
Clip & Expected & Predicted & Anomaly type & Score & \\
\hline
normal\_01           & NORMAL  & NORMAL  & none              & 0.00 & OK \\
normal\_02           & NORMAL  & NORMAL  & none              & 0.00 & OK \\
normal\_03           & NORMAL  & NORMAL  & none              & 0.00 & OK \\
missing\_cap\_01     & ANOMALY & ANOMALY & missing\_cap      & 0.90 & OK \\
missing\_cap\_02     & ANOMALY & ANOMALY & missing\_cap      & 0.95 & OK \\
wrong\_orient\_01    & ANOMALY & ANOMALY & wrong\_orientation& 0.90 & OK \\
wrong\_orient\_02    & ANOMALY & ANOMALY & wrong\_orientation& 0.90 & OK \\
wrong\_color\_01     & ANOMALY & ANOMALY & wrong\_color      & 1.00 & OK \\
missing\_pen\_01     & ANOMALY & ANOMALY & missing\_pen      & 1.00 & OK \\
motion\_stop\_01     & ANOMALY & ANOMALY & motion\_failure   & 0.90 & OK \\
motion\_stop\_02     & ANOMALY & ANOMALY & motion\_failure   & 0.90 & OK \\
motion\_temporal\_01 & ANOMALY & \textbf{NORMAL} & none      & 0.00 & \textbf{miss} \\
combined\_01         & ANOMALY & ANOMALY & wrong\_orientation& 0.90 & OK \\
\hline
\end{tabular}
\end{table}

\subsection{The Temporal Anomaly is a Perception Limit, Not a Wording Problem}
The unsolved clip (\texttt{temporal}: the pens shift on their own while the belt moves
normally) was missed by the v3 learner \emph{and} by the open-ended describe prompt,
which reported a perfectly normal scene. Because even an unconstrained description fails
to notice the motion, the cause is the VLM's sparse internal frame sampling, not the
wording of the questions. This motivated the alternative input representations below.

\subsection{Method 1 --- Native Video}
The default mode sends the clip as a video; DashScope samples frames internally. With the
v3 questions this reaches $0.92$ (single run) and catches the edge-orientation and the
belt-stop cases, but cannot see the intra-object motion.

\subsection{Method 2 --- Dense Frames (Central Window)}
Instead of relying on internal sampling, we extract $16$ time-ordered frames with OpenCV
and send them as separate numbered images, with a prompt focused on the
\emph{relative spacing} between adjacent pens (a cue that is invariant to the belt's
left-to-right translation). Sampling is restricted to the central window $[0.15, 0.70]$
of the clip to exclude the entry/exit frames. This representation \emph{catches} the
temporal anomaly ($0.9$ score) and keeps the normal clip normal, but, because the window
crop removes the edge frames and a late pause, it now misses the edge-orientation and the
belt-stop cases; overall accuracy is $0.85$.

\subsection{Method 3 --- Dense Frames (Wide Window)}
Widening the window to $[0.05, 0.95]$ to recover the lost cases backfired: the exit
frames, where the board carries the pens off-screen, were read as ``pens disappearing''
and produced a false positive on a normal clip. Accuracy fell to $0.62$ with one false
positive. This confirms the central crop was necessary, and that there is no free lunch.

\subsection{Method 4 --- Temporal Grid}
As a single-pass alternative we tiled the frames into one $3\times3$ montage image
(Figure~\ref{fig:grid}) so the model could compare moments side by side. On a four-clip
probe it caught the temporal case but shared the dense-frame blind spots (missed
orientation and stop) and, being full-span, re-introduced the exit false positive on the
normal clip. It therefore did not improve on the dense view.

\begin{figure}[!htbp]
    \centering
    \includegraphics[width=0.5\textwidth]{grid_temporal.jpg}
    \caption{\label{fig:grid}Temporal grid (3$\times$3 montage) of the temporal-anomaly
    clip. The three pens are evenly spaced early (cells 4--6) and then bunch together as
    one pen slides relative to the others; the relative-spacing change is what makes the
    intra-object motion visible.}
\end{figure}

\subsection{Method 5 --- Higher Video Sampling Rate (fps)}
We tested whether simply sampling the \emph{video} more densely would help, using the
\texttt{fps} parameter. Raising \texttt{fps} from the default to $8$ quadrupled the visual
tokens ($\sim$10{,}100 $\rightarrow$ $\sim$40{,}394), confirming denser sampling, yet the
full-question run still \emph{missed} the temporal anomaly and additionally lost the
edge-orientation case ($0.85$). Interestingly, a stand-alone call with $\texttt{fps}=8$
\emph{and} the dedicated relative-spacing prompt did catch the temporal anomaly while
keeping the normal clip normal --- showing that the decisive factor is the prompt focus
(an explicit relative-spacing question), not the frame count. Higher \texttt{fps} also
increases latency, so it is not a useful lever on its own.

\subsection{Method 6 --- Two-View Ensemble}
Since the video view and the dense-frame view have \emph{disjoint} blind spots and both
have zero false positives, we combine them with an OR rule: a clip is ANOMALY if either
view flags it. The video view supplies edge-orientation and belt-stop; the dense view
supplies the temporal anomaly. The union reaches $13/13$ with zero false positives at no
extra inference cost (it reuses the two existing runs).

\subsection{Representation Comparison}
Table~\ref{tab:repr} summarises the six methods. No single representation reaches
$13/13$: the continuous video captures the edge and stop cues but misses the intra-object
motion, whereas every discrete representation captures the motion but loses the other
cues. This \emph{continuous-versus-discrete} trade-off held across three discrete variants
(dense narrow, dense wide, grid). Only the ensemble closes the gap.

\begin{table}[!htbp]
\centering
\caption{\label{tab:repr}Input-representation comparison (blind, single run).}
\begin{tabular}{l|c|c|c|l}
\hline
Method & Accuracy & FP & Temporal? & Note \\
\hline
1. Video (v3, default fps)        & 0.92 (12/13) & 0 & no  & misses temporal \\
2. Dense [0.15--0.70], 16f        & 0.85 (11/13) & 0 & yes & misses orient-2, stop-2 \\
3. Dense wide [0.05--0.95]        & 0.62 (8/13)  & 1 & yes & exit artifact $\rightarrow$ FP \\
4. Temporal grid 3$\times$3       & ---$^{\dagger}$ & yes & yes & probe; FP + misses orient/stop \\
5. Video + fps=8 + v3             & 0.85 (11/13) & 0 & no  & fps degraded orientation \\
6. \textbf{Ensemble (1 $\vee$ 2)} & \textbf{1.00 (13/13)} & \textbf{0} & \textbf{yes} & union of blind spots \\
\hline
\end{tabular}
\\[2pt]
\footnotesize{$^{\dagger}$four-clip diagnostic probe, not a full $13$-clip run.}
\end{table}

\begin{figure}[!htbp]
    \centering
    \includegraphics[width=0.85\textwidth]{fig_accuracy.png}
    \caption{\label{fig:accuracy}Blind accuracy of each input representation on the
    $13$-clip pen dataset (single run). Only the ensemble reaches $1.00$.}
\end{figure}

\begin{figure}[!htbp]
    \centering
    \includegraphics[width=\textwidth]{fig_clip_matrix.png}
    \caption{\label{fig:matrix}Per-clip correctness of the two views and their ensemble.
    The video view misses only the temporal clip (\texttt{temp}); the dense-frame view
    misses only the edge-orientation (\texttt{or2}) and the late belt-stop (\texttt{st2}).
    Because the blind spots are disjoint and neither view produces false positives, the
    OR-ensemble is correct on all $13$ clips.}
\end{figure}

\subsection{Reasoning Budget and the Empty-Output Problem}
\texttt{qwen3.6-plus} is a reasoning model: on dense visual input it spent its entire
output budget on internal deliberation (e.g.\ $15{,}684$ reasoning tokens) and returned an
empty answer, crashing the run. Setting a \texttt{thinking\_budget} caps the reasoning and
forces a final JSON answer, which solved the empty-output problem and, as a side effect,
became our main latency lever. We also tried alternative models, but
\texttt{qwen-vl-max-latest} returned \emph{403 Access Denied} and \texttt{qwen3.7-max}
rejected video input, leaving \texttt{qwen3.6-plus} as the only usable model.
Figure~\ref{fig:reasoning} shows the per-clip reasoning tokens of the unbounded baseline:
a single hard clip consumed $13{,}394$ tokens, dominating that run's latency.

\begin{figure}[!htbp]
    \centering
    \includegraphics[width=0.85\textwidth]{fig_reasoning.png}
    \caption{\label{fig:reasoning}Reasoning tokens per clip under the unbounded baseline.
    The runaway outlier (clip \texttt{or2}) motivates capping the reasoning budget.}
\end{figure}

\subsection{Repeatability}
Because VLM outputs vary between runs, the baseline and the \texttt{thinking\_budget}=2000
configuration were each run three times under the blind protocol
(Table~\ref{tab:repeat}). In every run both configurations scored $11/13$ ($0.846$), so
the single-run $0.92$ reported earlier was a favourable draw and the true value is
$\approx 0.85$. The stable count, however, hides per-clip variability: the temporal clip
was missed in all six runs (a deterministic perception limit), but under the unbounded
baseline two borderline clips --- a normal clip and the edge-orientation clip --- changed
their NORMAL/ANOMALY verdict between runs, and the baseline even produced a false positive
in one run. Capping the reasoning budget removed this instability entirely: the three
\texttt{thinking\_budget}=2000 runs produced \emph{identical} per-clip verdicts (zero
flips). Reasoning-budget control therefore improves run-to-run determinism as well as
latency, and the result reframes every single-run number above as a noisy estimate.

\begin{table}[!htbp]
\centering
\caption{\label{tab:repeat}Repeatability over $n=3$ blind runs. ``Verdict flips'' is the
number of clips whose NORMAL/ANOMALY verdict was \emph{not} unanimous across the three
runs.}
\begin{tabular}{l|c|c|c|c}
\hline
Config & $n$ & Accuracy (each run) & Verdict flips & Total FP \\
\hline
baseline (unbounded reasoning) & 3 & $11/13$ $(0.846)$ & 2 clips & 1 \\
thinking\_budget = 2000        & 3 & $11/13$ $(0.846)$ & 0 clips & 0 \\
\hline
\end{tabular}
\end{table}

\begin{figure}[!htbp]
    \centering
    \includegraphics[width=0.85\textwidth]{fig_repeat.png}
    \caption{\label{fig:repeat}How often each clip was missed across the six repeated
    blind runs. The temporal clip (\texttt{temp}) is missed every time (a deterministic
    perception limit), whereas the borderline clips (\texttt{or2}, a normal clip) flip
    between runs.}
\end{figure}

\subsection{Latency vs.\ Reasoning Budget}
The reasoning generation, not the video prefill, dominates latency (one clip used
$13{,}394$ reasoning tokens). Capping it reduces latency at no accuracy cost
(Table~\ref{tab:latency}): \texttt{thinking\_budget}=2000 is about $41\%$ faster than the
unbounded baseline at the same mean accuracy, with one fewer false positive. Raising
\texttt{fps} instead makes the call slower and hurts accuracy, so it is the wrong lever
for speed.

\begin{table}[!htbp]
\centering
\caption{\label{tab:latency}Latency vs.\ reasoning budget (mean over $n=3$).}
\begin{tabular}{l|c|c}
\hline
Config & Seconds / clip & Accuracy \\
\hline
baseline (unbounded) & 85.2 & 0.846 \\
thinking\_budget = 2000 & 49.9 & 0.846 \\
thinking\_budget = 1000 & 39.5$^{*}$ & $\approx$0.85$^{*}$ \\
\hline
\end{tabular}
\\[2pt]
\footnotesize{$^{*}$single run; borderline clips more volatile.}
\end{table}

\begin{figure}[!htbp]
    \centering
    \includegraphics[width=0.7\textwidth]{fig_latency.png}
    \caption{\label{fig:latency}Mean latency per clip versus reasoning budget. Capping the
    budget gives a large speed-up at unchanged accuracy ($\approx 0.85$).}
\end{figure}

\subsection{Summary of the VERA-lite Methods}
The progression can be read as a sequence of motivated attempts: question optimisation
lifted accuracy to $0.92$; the one remaining failure was traced to a perception limit,
not wording; discrete representations recovered that failure but introduced others;
widening or re-tiling them traded one failure for a false positive; higher \texttt{fps}
showed presentation matters more than frame count; the ensemble combined the
complementary views to $13/13$; and finally, repeated runs reframed the honest accuracy as
$0.846 \pm 0.000$, with the reasoning budget giving a $41\%$ speed-up for free.
